\title{xLSTM: Extended Long Short-Term Memory}

\begin{abstract}
The introduction of the constant error carousel and gating mechanisms in the 1990s established Long Short-Term Memory (LSTM) as foundational concepts within the field. Over the years, LSTMs have demonstrated enduring utility, facilitating a variety of breakthroughs in deep learning—most notably, forming the basis of the earliest Large Language Models (LLMs). Nevertheless, the emergence of Transformers, characterized by their scalable, parallelizable self-attention, has shifted the paradigm and largely supplanted LSTMs for large-scale tasks. In this work, we reexamine language modeling performance by scaling up LSTMs to the multi-billion parameter regime, adapting current innovations from recent LLM research, while systematically addressing the well-documented drawbacks of conventional LSTMs. Our contributions are twofold. First, we present exponential gating, augmented with rigorous normalization and stabilization methods. Second, we re-envision the LSTM memory construct, leading to: (i) sLSTM, which features a scalar-based memory, scalar updates, and redesigned memory blending; and (ii) mLSTM, a highly parallelizable variant with matrix memory and a covariance-inspired update schema. By integrating these advanced LSTM designs as residual blocks, we form xLSTM modules, which are then stacked to assemble complete xLSTM networks. The adoption of exponential gating and these novel memory architectures enables xLSTM to achieve competitive—and at times superior—performance relative to contemporary Transformer architectures and State Space Models, both in terms of accuracy and scaling properties.\\
Code available at: \url{https://github.com/NX-AI/xlstm}
\end{abstract}

\section{Introduction}
\label{sec:intro}

The concept of Long Short-Term Memory (LSTM)~\citep{Hochreiter:91,Hochreiter:97nips,Hochreiter:97}, specifically the use of constant error carousel and gating mechanisms, was developed to address the vanishing gradient issue present in recurrent neural networks~\citep{Hochreiter:91,Hochreiter:00book}:

\begin{align}
\label{eq:lstm_idea}
  c_t \ = \ f_t \ c_{t-1} \ + \ 
          i_t \  z_t \ , \quad  
          h_t \ = \ o_t \ \psi({c_t}) \ . 
\end{align}

The constant error carousel functions through the additive modification of the cell state $c_{t-1}$ (depicted in green), utilizing cell inputs $z_t$ and regulated via sigmoid-activated gates (indicated in blue). The input gate $i_t$ together with the forget gate $f_t$ dictate the manner in which this update is executed. In contrast, the output gate $o_t$ determines the resultant output from the memory cell—namely, the hidden state $h_t$. Before this output is generated, the cell state undergoes normalization or compression using the function $\psi$, after which the output gating mechanism produces the final hidden state.

LSTMs have found widespread utility across a broad spectrum of fields \citep{Hochreiter:01,Hochreiter:07,Schmidhuber:15}, maintaining a dominant role in text generation until the rise of Transformers in 2017~\citep{Vaswani:17}. Their robust performance has been validated in a variety of sequence-based tasks, including but not limited to text synthesis~\citep{Graves:13,Karpathy:15blog}, handwritten content generation~\citep{Graves:13}, sequence-to-sequence translation~\citep{Sutskever:14nips}, automated program evaluation~\citep{Zaremba:14arxiv}, automated image description~\citep{Karpathy:15,Hossain:19}, source code production~\citep{Karpathy:15blog}, modeling rainfall-runoff dynamics~\citep{Kratzert:18,Kratzert:19}, and the development of hydrological models to issue flood warnings~\citep{Nearing:24}. Within reinforcement learning, LSTMs have consistently achieved top-tier performance in sequential modeling tasks; for instance, they were integral to notable systems such as AlphaStar for StarCraft II~\citep{Vinyals:17short}, OpenAI Five for Dota 2~\citep{OpenAI:19}, and models addressing magnetic control in nuclear fusion~\citep{Degrave:22short}. The capability of LSTMs to form and manipulate abstractions by proficiently extracting semantic features and retaining these in memory cells has been documented~\citep{Karpathy:15blog}. This property is illustrated by the emergence of number and syntax units~\citep{Lakretz:19}, linguistically specialized neurons~\citep{Bau:19}, and sentiment-sensitive cells~\citep{Radford:17arxiv}. LSTMs remain integral to many contemporary, high-impact applications~\citep{Degrave:22short,Nearing:24}, confirming their enduring significance.

While LSTMs have achieved remarkable results, they are hindered by three primary drawbacks: (i) They lack the capability to adjust previously made storage decisions. This deficiency is highlighted using the {\it Nearest Neighbor Search} problem (refer also to Appendix~\ref{sec:appNearestNeighborSearch}): when provided with a reference vector, the task requires scanning through a sequence to locate the most similar vector and returning its corresponding value at the end of the sequence. The left panel of Figure~\ref{fig:lstmProblems} depicts the mean squared error obtained for this task. LSTMs face difficulty in updating a stored value when a more similar vector appears later in the sequence, whereas our novel xLSTM addresses this shortfall through the use of exponential gating. (ii) LSTMs possess restricted storage capabilities, requiring all information to be condensed into scalar cell states. This challenge is demonstrated by the {\it Rare Token Prediction} task. As illustrated in the right panel of Figure~\ref{fig:lstmProblems}, perplexity scores on Wikitext-103~\citep{Merity:17} are reported for tokens grouped by frequency. Due to constrained storage, LSTMs underperform with rare tokens. This issue is mitigated in our new xLSTM model, which employs matrix-based memory. (iii) Their inherent memory mixing through hidden-to-hidden connections across consecutive time steps prevents parallel computation and enforces sequential data processing.

The shortcomings of LSTM have contributed to the rise of Transformers~\citep{Vaswani:17} within language modeling tasks. This raises the question: by addressing these issues and expanding LSTM architectures to match the scale of contemporary Large Language Models, what levels of performance might be attainable in language modeling?

\section{Extended Long Short-Term Memory}
\label{sec:xLSTM}

In addressing the constraints inherent in LSTM, the Extended Long Short-Term Memory (xLSTM) architecture implements two significant enhancements to the conceptual formulation of Equation~\eqref{eq:lstm_idea}. These enhancements—namely, exponential gating and innovative memory structures—broaden the LSTM family by introducing two variants: (i) sLSTM (see Section~\ref{sec:sLSTM}), characterized by a scalar memory, scalar updates, and mechanisms for memory mixing, and (ii) mLSTM (see Section~\ref{sec:mLSTM}), which utilizes a matrix memory and employs a covariance (outer product) update rule that facilitates complete parallelization. The adoption of exponential gating serves as a common augmentation in both sLSTM and mLSTM. To realize efficient parallel execution, mLSTM forgoes the traditional memory mixing, effectively removing hidden-to-hidden recurrent connections. Both architectures support generalizations to multiple memory cells; in the sLSTM case, memory mixing occurs across these cells. The sLSTM can additionally incorporate multiple heads—without mixing memory across heads but preserving intra-head memory mixing among cells—thereby enabling a novel form of memory mixing afforded by the combination of heads and exponential gating. For mLSTM, the distinction between multiple heads and multiple cells does not hold, as the two concepts are interchangeable.

By incorporating these advanced LSTM variants within residual block structures, we obtain xLSTM blocks (refer to Section~\ref{sec:xLSTMblock}). Sequentially layering these xLSTM blocks using residual connections yields what we term xLSTM architectures (see Section~\ref{sec:xLSTMarchitecure}).
Figure~\ref{fig:xlstm_sketch} provides an illustration of the xLSTM architecture and its constituent elements.

\subsection{Review of the Long Short-Term Memory}
\label{sec:orgLSTM}

The foundational concept behind LSTM~\citep{Hochreiter:91,Hochreiter:97nips,Hochreiter:97} featured a scalar memory cell that serves as both the computational and memory core, effectively circumventing the vanishing gradient problem~\citep{Hochreiter:91,Hochreiter:00book} via a mechanism termed the constant error carousel, represented by the cell state update. The architecture of the memory cell incorporates three gates: input gate, output gate, and forget gate. The forget gate was first proposed by \citet{Gers:00}. The memory cell’s update equations at time step $t$ are as follows:

\begin{align}
c_t \ &= \  f_t \ c_{t-1} \ + \ i_t \ z_t &  & &\text{cell state} \\
h_t  \ &= \ o_t \ \tilde{h}_t \ , & \tilde{h}_t \ &= \ \psi \left( c_t\right)
  &\text{hidden state} \\
z_t \ &= \ \varphi \left( \tilde{z}_t \right) \ , 
  &\tilde{z}_t \ &=  \ \mathbf{w}^\top_{z} \ \mathbf{x}_t \ + \
  r_{z}  h_{t-1} \ + \  b_{z} \ \
  &\text{cell input} \\
i_t \ &= \ \sigma \left( \tilde{i}_t  \right) \ , 
  &\tilde{i}_t \ &= \ \mathbf{w}^\top_{i} \ \mathbf{x}_t \ + \
  r_{i}  \ h_{t-1} \ + \  b_{i} \ \
  &\text{input gate} \\
f_t \ &= \ \sigma \left(  \tilde{f}_t \right) \ , 
  &\tilde{f}_t \ &= \ \mathbf{w}^\top_{f} \ \mathbf{x}_t  \ + \
  r_{f}  \ h_{t-1} \ + \  b_{f} \ \
  &\text{forget gate} \\
o_t \ &= \ \sigma \left( \tilde{o}_t \right) \ , 
  &\tilde{o}_t  \ &= \ \mathbf{w}^\top_{o} \ \mathbf{x}_t \ + \
  r_{o}  \ h_{t-1} \ + \  b_{o} \ \
  &\text{output gate} 
\end{align}

The vectors $\mathbf{w}_{z}$, $\mathbf{w}_{i}$,
$\mathbf{w}_{f}$, and $\mathbf{w}_{o}$ denote the input weights connecting the input $\mathbf{x}_t$ to the cell input, input gate, forget gate, and output gate, respectively. The parameters $r_{z}$, $r_{i}$, $r_{f}$, and $r_{o}$ are the recurrent weights linking the previous hidden state $h_{t-1}$ to the cell input, input gate, forget gate, and output gate, respectively. The terms $b_{z}$, $b_{i}$, $b_{f}$, and $b_{o}$ represent the corresponding biases. The activation functions $\varphi$ and $\psi$ govern the cell input and the hidden state (often chosen as $\tanh$); the function $\psi$ in particular prevents the cell state from becoming unbounded by squashing or normalizing its values. For all gating units, the sigmoid function serves as the activation, that is, $\sigma \left( x \right)= 1/(1+\exp(-x))$. Later models aggregate multiple scalar memory states $\mathbf{c}_t \in \mathbb{R}^{d}$ into a vector, enabling the deployment of recurrent weight matrices $\mathbf{R} \in \mathbb{R}^{d \times d}$ which facilitate interactions among different memory cells’ outputs~\citep{Greff:15}; further explanation is given in Appendix~\ref{sec:apporgLSTM}. Comprehensive ablation analyses indicate that each element of the memory cell contributes essentially to its performance~\citep{Greff:15}.

\subsection{sLSTM}
\label{sec:sLSTM}

To enhance LSTMs so that they can reconsider their storage choices, we incorporate exponential gates (red) as well as mechanisms for normalization and stabilization. Specifically, exponential activation functions are utilized in the input and forget gates. For normalization purposes, a dedicated normalizer state is added, which accumulates the sum of products between the input gate and all subsequent forget gates.

The scalar sLSTM forward pass is:

\begin{align}
\label{eq:slstmforward}
c_t \ &= \  f_t \ c_{t-1} \ + \ i_t \ z_t & & 
 &\text{cell state} \\
n_t \ &= \  f_t \ n_{t-1} \ + \ i_t  & & 
 &\text{normalizer state} \\
h_t \ &= \ o_t \ \tilde{h}_t \ ,
  &\tilde{h}_t \ &= \ c_t / n_t
&\text{hidden state} \\
z_t \ &= \ \varphi \left( \tilde{z}_t \right) \ , 
  &\tilde{z}_t \ &=  \ \mathbf{w}^\top_{z} \ \mathbf{x}_t \ + \
  r_{z}  h_{t-1} \ + \  b_{z}
  &\text{cell input} \\
i_t \ &= \ \!\exp \left( \tilde{i}_t  \right) \ , 
  &\tilde{i}_t \ &= \ \mathbf{w}^\top_{i} \ \mathbf{x}_t \ + \
  r_{i}  \ h_{t-1} \ + \  b_{i}
  &\text{input gate} \\
f_t \ &= \ \sigma \left(  \tilde{f}_t \right) \ \text{OR} \
\!\exp \left(  \tilde{f}_t \right) \ ,
  &\tilde{f}_t \ &= \ \mathbf{w}^\top_{f} \ \mathbf{x}_t  \ + \
  r_{f}  \ h_{t-1} \ + \  b_{f}
  &\text{forget gate} \\
o_t \ &= \ \sigma \left( \tilde{o}_t \right) \ , 
  &\tilde{o}_t  \ &= \ \mathbf{w}^\top_{o} \ \mathbf{x}_t \ + \
  r_{o}  \ h_{t-1} \ + \  b_{o}
  &\text{output gate} 
\end{align}

The conventional gating mechanisms from LSTM models—namely, input and/or hidden state-dependent gates along with a bias component—are adapted for use in the proposed architectures. Since exponential activation functions have the potential to generate excessively large outputs, which may result in overflow, we mitigate this issue by introducing an auxiliary state \mbox{$m_t$~\citep{Milakov:18arxiv}} to enhance gate stability:

\begin{align}
\label{eq:slstmstabil}
m_t \ &= \ \max \left( \log ( f_t ) + m_{t-1} , \log ( i_t ) \right) &\text{stabilizer state} \\
i'_t \ &= \ \!\exp \left( \log \left ( i_t \right) - m_t \right) \ = \ \!\exp \left( \tilde{i}_t - m_t \right) \ \, 
  &\text{stabil. input gate} \\
f'_t \ &= \ \!\exp \left( \log \left( f_t \right) + m_{t-1} - m_t \right) \ \, 
  &\text{stabil. forget gate}
\end{align}

As demonstrated in Appendix~\ref{sec:appsLSTM}, substituting $f_t$ with $f'_t$ and $i_t$ with $i'_t$ during the forward computation has no effect on either the overall network output or on the gradients of the loss relative to the model parameters.

\paragraph{New Memory Mixing.} Similar to the original LSTM architecture (refer to Appendix~\ref{sec:appsLSTM}), sLSTM supports several memory cells. The presence of multiple memory cells allows for mixing of information through recurrent pathways $\mathbf{R}_{\mathbf{z}}$, $\mathbf{R}_{\mathbf{i}}$, $\mathbf{R}_{\mathbf{f}}$, and $\mathbf{R}_{\mathbf{o}}$ that link the hidden state $\mathbf{h}$ to the memory cell input $\mathbf{z}$ as well as to the gating variables $\mathbf{i}$, $\mathbf{f}$, and $\mathbf{o}$. What differentiates the memory mixing in this context is the incorporation of exponential gating. In this revised sLSTM, it is possible to employ several heads, each enabling memory mixing solely within itself—without inter-head interactions. By combining the head-based structure of sLSTM with exponential gating, a novel mechanism for mixing memory representations is achieved.

\subsection{mLSTM}
\label{sec:mLSTM}

To improve the storage capability of LSTMs, the traditional scalar memory cell $c \in \mathbb{R}$ is generalized to a matrix form $\mathbf{C} \in \mathbb{R}^{d \times d}$. This modification enables the retrieval process through matrix multiplication. Specifically, at time step $t$, we store two vectors: a key $\mathbf{k}_t \in \mathbb{R}^d$ and an associated value $\mathbf{v}_t \in \mathbb{R}^d$ (adopting terminology from the Transformer architecture). Subsequently, at time $t + \tau$, the value $\mathbf{v}_t$ is to be accessed by means of a query vector $\mathbf{q}_{t+\tau} \in \mathbb{R}^d$. This formulation aligns with the paradigm of Bidirectional Associative Memories (BAMs) as established by prior work \citep{Kohonen:72,Anderson:72,Nakano:72,Anderson:77}.

The covariance update rule~\citep{Sejnowski:77,Dayan:91} for storing a key-value pair is

\begin{align}
\mathbf{C}_t \ &= \ \mathbf{C}_{t-1} \ + \ \mathbf{v}_{t} \ \mathbf{k}_t^\top \ .
\end{align}

We posit the application of layer normalization prior to mapping inputs into key and value representations, ensuring these projections maintain a zero-mean property. The covariance update mechanism, as demonstrated to be optimal~\citep{Dayan:91}, facilitates maximal distinguishability among retrieved binary vectors, correlating to an optimal signal-to-noise ratio. When considering solely pairwise retrieval interactions and accepting the quadratic computational overhead akin to attention~\citep{Krotov:16,Krotov:17,Ramsauer:21}, even greater separability can be obtained. Furthermore, the covariance update rule aligns with the methodology of Fast Weight Programmers~\citep{Schmidhuber:92ncfastweights,Schlag:21}, which have subsequently incorporated a constant decay factor applied to $\mathbf{C}_{t-1}$ as well as a fixed learning rate assigned to 
$\mathbf{v}_{t} \mathbf{k}_t ^\top$~\citep{Ba:16}.
Following this perspective, we embed the covariance update approach within the LSTM architecture, interpreting the forget gate as the decay factor, the input gate as the learning rate, and allowing the output gate to modulate the retrieved vector.

In this matrix memory architecture, the normalizer state is computed as a weighted aggregation of key vectors, with each key vector's contribution determined by the product of the input gate and subsequent forget gates. This ensures that the normalizer state accurately reflects the accumulated gating strengths. Given that the dot product between the query and the normalizer state may approach zero, its absolute value is taken and then constrained to a minimum threshold (usually set to 1.0), in accordance with the approach described in~\citep{Sun:23arxiv}. The forward computation for mLSTM is as follows:

\begin{align}
\label{eq:mlstm_recurrent_begin}
\mathbf{C}_t \ &= \ f_t \, \mathbf{C}_{t-1} \ + \ 
  i_t \, \mathbf{v}_t \, \mathbf{k}_t^\top & & &\text{cell state} \\
\mathbf{n}_t \ &= \ f_t \, \mathbf{n}_{t-1} \ + \
  i_t \, \mathbf{k}_t & & &\text{normalizer state} \\
\mathbf{h}_t &= \ \mathbf{o}_t \odot \tilde{\mathbf{h}}_t
\ , \qquad \qquad 
\tilde{\mathbf{h}}_t \ = \ \dfrac{\mathbf{C}_t \mathbf{q}_t}{
  \max \left\{ |\mathbf{n}_t^\top \mathbf{q}_t|,\, 1 \right\} }
  &&&\text{hidden state} \\
\mathbf{q}_t \ &= \ \mathbf{W}_q \mathbf{x}_t + \mathbf{b}_q & & &\text{query input} \\
\mathbf{k}_t \ &= \ \frac{1}{\sqrt{d}} \mathbf{W}_k \mathbf{x}_t + \mathbf{b}_k & & &\text{key input} \\
\mathbf{v}_t \ &= \ \mathbf{W}_v \mathbf{x}_t + \mathbf{b}_v & & &\text{value input} \\
i_t \ &= \exp\!\left(\tilde{i}_t\right),  \qquad \quad \tilde{i}_t = \mathbf{w}_i^\top \mathbf{x}_t + b_i
  &&&\text{input gate} \\
f_t \ &= \sigma(\tilde{f}_t) \ \text{OR} \ \exp(\tilde{f}_t),\ \ 
\tilde{f}_t = \mathbf{w}_f^\top \mathbf{x}_t + b_f
  &&&\text{forget gate} \\
\label{eq:mlstm_recurrent_end}
\mathbf{o}_t &= \sigma(\tilde{\mathbf{o}}_t), \qquad \quad \tilde{\mathbf{o}}_t = \mathbf{W}_{\mathbf{o}} \mathbf{x}_t + \mathbf{b}_{\mathbf{o}}
  &&&\text{output gate} 
\end{align}

The architecture of mLSTM supports several memory cells, similar to the original LSTM design. In mLSTM, utilizing multiple heads yields the same result as employing multiple cells, because memory contents are not shared or mixed across these components. To ensure stable behavior of the exponential gating mechanism in mLSTM, we employ stabilization methods identical to those adopted in sLSTM (refer to Equation~\ref{eq:slstmstabil}). Due to the absence of memory mixing in mLSTM, the recurrence relations admit a reformulation that allows for parallel computation. Additional information can be found in Appendix~\ref{sec:appmLSTM}.

\subsection{xLSTM Architecture}
\label{sec:xLSTMarch}

\paragraph{xLSTM Blocks.}
\label{sec:xLSTMblock}
An xLSTM block is intended to encapsulate past information through non-linear operations in a high-dimensional representation, thereby facilitating the distinction between various preceding contexts or histories. Such discrimination among different histories is essential for making accurate predictions about subsequent elements in a sequence, for example, the next token. This design philosophy draws inspiration from Cover’s Theorem~\citep{Cover:65}, which suggests that patterns transformed non-linearly into higher-dimensional spaces are more amenable to linear separation than those in their original representation. 

We examine two types of residual block architectures: (i) The first architecture is a residual block where the up-projection occurs after non-linear summarization (similar to the approach in Transformers). In this scheme, the model initially produces a non-linear summary of historical information in the input space, then applies a linear projection to elevate this representation to a higher-dimensional space, follows this with a non-linear activation, and finally projects it linearly back to the input space. This configuration is illustrated in the left panel of Figure~\ref{fig:Backbones} and the third column of Figure~\ref{fig:xlstm_sketch}. A comprehensive schematic is available in Figure~\ref{fig:appsLSTM_detailed} within the appendix. (ii) The second architecture implements the up-projection before non-linearity (reminiscent of State Space Models), where the input is first linearly mapped into a higher-dimensional space,

The approach involves a nonlinear aggregation of historical information within a high-dimensional space, followed by a linear projection back into the original feature space. When deploying an xLSTM block equipped with an sLSTM, the up-projection block is applied after processing. Conversely, for an xLSTM block with an mLSTM, the up-projection occurs prior to processing, as the memory capabilities are expanded in the high-dimensional context. Additional clarification is available in the left panel of Figure~\ref{fig:Backbones}, the third column of Figure~\ref{fig:xlstm_sketch}, and Figure~\ref{fig:appsLSTM_detailed} in the appendix.

\paragraph{xLSTM Architecture. }
\label{sec:xLSTMarchitecure}
The xLSTM model is designed by stacking modular blocks in a residual fashion~\citep{Srivastava:15,He:16}. Our implementation follows the widely adopted pre-LayerNorm~\citep{Ba:16b} residual backbone architecture, aligning with the approach found in modern Large Language Models. Refer to the final column of Figure~\ref{fig:xlstm_sketch} for an illustration.

\subsection{Memory and Speed Considerations}

In contrast with Transformers, xLSTM architectures exhibit linear computational requirements and fixed memory usage regardless of the input sequence length. Due to the compressive nature of xLSTM memory, these networks are particularly advantageous for use in industrial scenarios and for deployment in edge environments.

The mLSTM’s memory operates without needing additional parameters, yet it incurs substantial computational cost due to its $d \times d$ memory matrix and corresponding $d \times d$ update operations. This approach balances the scale of memory capacity with the demands of computational resources. However, since these operations are highly parallelizable on GPUs, their impact on overall runtime remains minimal.

Although mLSTM allows for parallelization similarly to FlashAttention~\citep{Dao:22,Dao:23} or GLA~\citep{Yang:23arxiv}, this is not the case for sLSTM because of the interdependence created by memory mixing (hidden-hidden connections). Nevertheless, we introduced an efficient CUDA implementation with register-level GPU memory optimizations, resulting in sLSTM running at no more than twice the computational time of mLSTM.

\section{Related Work}
\label{sec:related}

\paragraph{Linear Attention.} A variety of approaches have been proposed to address the quadratic scaling of Transformer attention relative to sequence length and to achieve linear complexity. The Synthesizer circumvents token-to-token dependencies by learning artificial attention weights~\citep{Tay:20arxiv}. Linformer employs a low-rank decomposition to implement self-attention and attains a linear approximation of this mechanism~\citep{Wang:20arxiv}. Linear Transformer modifies the computation of attention into a linear process \citep{Katharopoulos:20}. Performer employs positive orthogonal random features to approximate the attention softmax in a linear fashion~\citep{Choromanski:21}. Furthermore, fast, long-range convolutions have been utilized to supplant attention in architectures such as Structured Global Convolution (SGConv)~\citep{Li:22} and the Hyena Hierarchy~\citep{Poli:23}.

\paragraph{State Space Models.} In recent times, State Space Models (SSMs) have gained significant attention due to their linear computational complexity with respect to sequence length, as well as their competitive results in comparison to Transformer architectures. The early foundational work in this area was the Structured State Space sequence model (S4)~\citep{Gu:21}. Since then, several other architectures have emerged, including the Diagonal State Space (DSS) model~\citep{Gupta:22}, Gated State Space (GSS) models~\citep{Mehta:22}, S5 model~\citep{Smith:22}, Bidirectional Gated SSM (BiGS)~\citep{Wang:22}, H3 model~\citep{Fu:23}, and Mamba~\citep{Gu:24arxiv}.

\paragraph{Recurrent Neural Networks.} Recurrent Neural Networks (RNNs) have emerged as alternatives to Transformer-based models and attention mechanisms, primarily because their computation scales linearly with input sequence length. Notably, RNN variants equipped with Deep Linear Recurrent Units (LRUs) have achieved strong performance in language modeling tasks~\citep{Orvieto:23, De:24arxiv}. Similarly, architectures such as Hierarchically Gated Linear RNN (HGRN)~\citep{Qin:23} and its successor HGRN2~\citep{Qin:24arxiv} have also reported favorable outcomes. Among RNN methods for large language models, the RWKV approach~\citep{Peng:23arxivshort,Peng:24arxivshort} is particularly recognized for its performance that rivals that of Transformer models.

\paragraph{Gating.}
Gating, originally introduced as a central concept in LSTM, has since been revived and reconceptualized by numerous modern methods. Models such as HGRN~\citep{Qin:23}, HGRN2~\citep{Qin:24arxiv}, Gated Linear Attention (GLA)~\citep{Yang:23arxiv}, Gated State Space (GSS) models~\citep{Mehta:22}, Bidirectional Gated SSM (BiGS)~\citep{Wang:22}, Moving Average Equipped Gated Attention (MEGA)~\citep{Ma:22}, RWKV~\citep{Peng:23arxivshort}, and Mamba~\citep{Gu:24arxiv} all incorporate gating mechanisms.

\paragraph{Covariance Update Rule.} To improve memory capacity, we integrated a matrix memory into the mLSTM cell, leveraging a covariance-based update rule. Several other architectures have adopted similar mechanisms, including Fast Weight Programmers \citep{Schmidhuber:92ncfastweights,Schlag:21}, RWKV-5 and RWKV-6~\citep{Peng:24arxivshort}, Retention~\citep{Sun:23arxiv}, Linear Transformer~\citep{Katharopoulos:20}, and HGRN2~\citep{Qin:24arxiv}.

\paragraph{Most Related.} Among existing architectures, the designs most analogous to xLSTM are Retention~\citep{Sun:23arxiv}, RWKV~\citep{Peng:23arxivshort,Peng:24arxivshort}, and HGRN2~\citep{Qin:24arxiv}. All of these systems incorporate elements such as matrix memory and gating mechanisms. Nevertheless, unlike the recently proposed sLSTM, these models lack the mechanism for memory mixing. The incorporation of memory mixing allows for effective state tracking, rendering LSTMs inherently more expressive than both State Space Models (SSMs) and Transformers~\citep{Merrill:24,Deletang:23}. Robust state tracking is essential for tasks such as code evaluation or following entities throughout extended narratives.

\paragraph{Residually Stacking Architectures.} Similar to the majority of state-of-the-art deep learning frameworks, xLSTM architectures are organized by stacking basic modules in a residual fashion~\citep{Srivastava:15,He:16}.

This architectural paradigm facilitated the advancement of deep convolutional networks~\citep{He:16} as well as the development of Transformers~\citep{Vaswani:17}. Transformers, in turn, have served as the foundational engine for Large Language Models (LLMs), including models such as GPT-3~\citep{Brown:20short}, ChatGPT~\citep{Schulman:22short}, GPT-4~\citep{Achiam:23gpt4short}, Megatron-LM~\citep{Shoeybi:19}, Gopher~\citep{Rae:21short}, ERNIE 3.0 Titan~\citep{Wang:21arxivshort}, GLaM~\citep{Du:21glamshort}, Chinese M6~\citep{Lin:21}, multilingual AlexaTM 20B~\citep{Soltan:22}, OPT~\citep{Zhang:22opt}, Chinchilla~\citep{Hoffmann:22short}, BLOOM~\citep{Scao:22arxivshort}, GLM-130B~\citep{Zeng:22arxivshort}, LaMDA~\citep{Thoppilan:22short}, PaLM~\citep{Chowdhery:22short}, Llama~\citep{Touvron:23arxiv}, and Gemini~\citep{GeminiTeam:23arxiv,Reid:24arxiv}.

\section{Experiments}
\label{sec:Exp}

In this section, we present an empirical analysis of xLSTM, positioning it alongside established techniques with an emphasis on language modeling applications. Section~\ref{sec:ExpSyntethic} delves into xLSTM's unique strengths using synthetic benchmarks. In Section~\ref{sec:ExpComparison}, we benchmark the validation perplexity for a selection of prominent language modeling architectures, all trained on a 15B token split from SlimPajama~\citep{Soboleva:23}, and additionally conduct ablation experiments specifically for xLSTM. We further analyze scaling properties in line with the methodologies of \citet{Kaplan:20} and \citet{Brown:20short}. 

Moving to Section~\ref{sec:ExpLanguage}, we extend our evaluation to a comprehensive language modeling study. Here, we pit xLSTM against the top contenders identified in Section~\ref{sec:ExpComparison}, with each method being trained on a substantially larger corpus of 300B tokens from SlimPajama~\citep{Soboleva:23}. Our analysis proceeds in several parts: (1) we measure extrapolation abilities to handle extended context lengths, (2) we examine validation perplexity and assess effectiveness on downstream tasks as per~\citep{Sutawika:24}, (3) we evaluate generalization across 571 distinct textual domains from the PALOMA language benchmark~\citep{Magnusson:23arxivshort}, and (4) we revisit scaling analysis for all approaches using a dataset twenty times the original training size.

\label{sec:xlstm_blocks_notation}
Throughout all experiments, we denote xLSTM[$a$:$b$] to represent the proportion $a/b$ of xLSTM blocks built with mLSTM cells versus those constructed with sLSTM cells. As an illustration, xLSTM[7:1] indicates a configuration where, among eight blocks, seven employ mLSTM cells and one incorporates an sLSTM cell. In scenarios where the total number of blocks is 48, this allocation corresponds to 42 mLSTM-based blocks and 6 sLSTM-based blocks. Additionally, in every experiment, mLSTM blocks utilize pre up-projection, while sLSTM blocks apply post up-projection mechanisms.

\subsection{Synthetic Tasks and Long Range Arena}
\label{sec:ExpSyntethic}
To begin, we evaluate xLSTM’s novel exponential gating with memory mixing mechanisms using formal language benchmarks~\citep{Deletang:23}. Subsequently, we examine the performance of xLSTM’s new matrix memory component in the Multi-Query Associative Recall task~\citep{Arora:23arxiv}. Lastly, we investigate xLSTM’s capability for handling lengthy input sequences by conducting experiments on the Long Range Arena benchmark~\citep{Tay:21}.

\paragraph{Evaluation of Exponential Gating with Memory Mixing in xLSTM.} We examine the capabilities of xLSTM's newly introduced exponential gating mechanism combined with memory mixing, which is intended to facilitate the resolution of state tracking challenges~\citep{Merrill:24, Merrill:23}. Building on and extending the suite of formal language experiments from \citet{Deletang:23}, we modify the tasks to support multi-length training as required for assessing extrapolation to longer sequences. Comprehensive descriptions of the tasks and supplementary results can be found in Appendix~\ref{sec:appExpSynthetic-formalLang}. Our assessment involves benchmarking xLSTM against a range of models, including Transformers, State Space Models, and traditional Recurrent Neural Networks. Model performance is measured in terms of accuracy on those tokens essential to each specific task, with accuracy normalized between 0 (equivalent to random guessing) and 1 (flawless performance). The study involves comparing two-block variants of the models, specifically: xLSTM[0:1] (sLSTM only), xLSTM[1:0] (mLSTM only), xLSTM[1:1], along with Llama, Mamba, RWKV, Retention, Hyena, LSTM, and Transformer-embedded LSTM blocks (LSTM (Block)). The experimental outcomes are summarized in Figure~\ref{fig:formal-main}. We observe that models like Transformers and State Space Models lacking memory mixing mechanisms—and thus, the ability to maintain state—are unable to resolve certain regular language tasks, such as the parity problem. These findings corroborate previous research indicating that Transformers and State Space Models are inherently less capable than Recurrent Neural Networks with respect to state tracking~\citep{Merrill:24, Merrill:23, Deletang:23}.

\paragraph{Test of xLSTM's Memory Capacities on Associative Recall Tasks.} This study assesses the newly introduced matrix memory of xLSTM by examining its memory performance on the Multi-Query Associative Recall task~\citep{Arora:23arxiv}. In this setup, sequences consist of key-value pairs, selected at random from an extensive vocabulary, which must be retained for subsequent retrieval. To further challenge the baseline task, we increase the quantity of key-value pairs to as many as 256, and the context length is expanded up to 2048 tokens. This design allows for a more rigorous evaluation of the memory limits across different model architectures. The analysis involves 2-block variants of Llama, Mamba, RWKV-5, RWKV-6, xLSTM[1:1], and xLSTM[1:0], comparing their accuracy in retrieving stored pairs. Given that Transformers such as Llama feature memory that scales exponentially with the coding dimension~\citep{Ramsauer:21}, they are considered the benchmark for this assessment. The outcomes, presented in Figure~\ref{fig:mqar-main}, reveal that xLSTM[1:1] outperforms all non-Transformer architectures, including in lower capacity models. Notably, incorporating the sLSTM block does not compromise memory ability; rather, it enhances capacity, especially evident in tasks involving 256 key-value pairs, the most demanding scenario evaluated. Further evidence, including extrapolation experiments, is provided in Appendix~\ref{sec:appExpSynthetic-mqar}, demonstrating that xLSTM's advanced memory also supports generalization to longer contexts not observed during the training phase.

\paragraph{Test of xLSTM's Long Context Capabilities on Long Range Arena.} In order to evaluate how effectively xLSTM manages extended sequences and substantial context windows, we conduct a comparative analysis with other approaches on the Long Range Arena~\citep{Tay:21}. The results indicate that xLSTM consistently achieves high performance across the range of tasks, implying that its architecture is highly adept at tackling various challenges associated with long context scenarios.

For more details, see Appendix~\ref{sec:appExpSynthetic-lra}.

\subsection{Method Comparison and Ablation Study}
\label{sec:ExpComparison}

In order to investigate the central research question—specifically, the performance of our proposed LSTM variants when scaled for language modeling—we conduct training of xLSTMs, Transformers, State Space Models, and additional techniques using 15B tokens sourced from SlimPajama, maintaining a consistent auto-regressive framework. We assess all trained models on the validation split and carry out ablation analyses focused on the xLSTM architectures.

\paragraph{Comparing xLSTM to Other Methods.} In our comparative analysis, models are trained on a dataset of 15 billion tokens derived from SlimPajama~\citep{Soboleva:23}. The assessment metric for the models is perplexity measured on a held-out validation set. We benchmark xLSTM (our proposed approach) against a variety of established techniques: GPT-3 (Transformer)~\citep{Brown:20short}, Llama (Transformer)~\citep{Touvron:23arxiv}, H3 (SSM)~\citep{Fu:23}, Mamba (SSM)~\citep{Gu:24arxiv}, RWKV-4 (RNN)~\citep{Peng:23arxivshort}, RWKV-5 (RNN)~\citep{Peng:24arxivshort}, RWKV-6 (RNN)~\citep{Peng:24arxivshort}, GLA (linear Transformer)~\citep{Yang:23arxiv}, HGRN (RNN)~\citep{Qin:23}, HGRN2 (RNN)~\citep{Qin:24arxiv}, RetNet (linear Transformer)~\citep{Sun:23arxiv}, Hyena (linear Transformer)~\citep{Poli:23}, as well as two xLSTM variants, xLSTM[1:0] and xLSTM[7:1] (for further details, refer to Section~\ref{sec:xlstm_blocks_notation}). Training procedures utilize mixed-precision; however, for RWKV-5, RWKV-6, GLA, and HGRN2, PyTorch's automated mixed precision is not adopted (refer also to Appendix Section~\ref{sec:appGenTrainProc} for additional training specifics).

The evaluated models are grouped into the following categories: (a) Transformers, (b) State Space Models (SSMs), and (c) both Recurrent Neural Networks (RNNs) and linear Transformer variants—where the latter are characterized by replacing the canonical attention mechanism with a linear approximation. To ensure a fair comparison, all models are configured with parameter count parity to a 350M-parameter GPT-3 baseline, employing an embedding dimension of 1024 and 24 residual blocks. GPT-3 distinguishes itself in the parameter count by virtue of tying its token and output embedding weights, thus resulting in a modest reduction in overall parameters. According to results reported in Table~\ref{tab:spaj15b_model_results}, xLSTM achieves the lowest validation perplexity among all compared approaches. Further experimental details are available in Appendix~\ref{sec:appExpComparison}. Figure~\ref{fig:temp_sclaw15B} illustrates how these models scale in performance, demonstrating that xLSTM is expected to maintain or surpass current benchmarks as model size increases.

\paragraph{Ablation Studies.}
As shown in Table~\ref{tab:spaj15b_model_results} and Figure~\ref{fig:temp_sclaw15B}, xLSTM attains strong language modeling performance after being trained on 15B tokens from SlimPajama. To dissect the individual contributions leading from LSTM to xLSTM, we incrementally transform a standard LSTM model into the xLSTM variant. The process begins by embedding LSTM layers into residual backbones with pre-LayerNorm. Next, we incorporate a block with post up-projection. The final modifications involve adding both exponential gating mechanisms and matrix memory.

The results are shown in Table~\ref{tab:ablstudies} (top). 

The notable gains in performance highlighted by the ablation studies stem from both the matrix memory and the exponential gating mechanisms. Recognizing the significance of gating in RNNs and State Space Models, we also investigate the impact of various gating strategies. Table~\ref{tab:ablstudies} (bottom) demonstrates that allowing each gate to be learned and modulated by the input produces a cumulative benefit. Further analyses concerning the distinct backbone components are provided in Appendix~\ref{sec:appAblDetails}.

\subsection{xLSTM as Large Language Model}
\label{sec:ExpLanguage}

This investigation concludes with comprehensive language modeling experiments aimed at evaluating xLSTM’s viability as a large language model (LLM). For this purpose, we scale up the training dataset, utilizing 300B tokens derived from SlimPajama. Notably, this token count matches the training regimens employed by models such as Mamba~\citep{Gu:24arxiv} and Griffin~\citep{De:24arxiv}.

We benchmark xLSTM against RWKV-4, Llama, and Mamba—selecting a representative from each methodological category discussed in Section~\ref{sec:ExpComparison}. RWKV-4 is used to represent RNN-based models, as viable training precision parameters for RWKV-5, RWKV-6, and HGRN2 were only established after initiation of the 300B token training experiments (refer to Appendix~\ref{sec:appGenTrainProc}). For each method, we train several parameter scales (125M, 350M, 760M, 1.3B), probe all configurations for their sequence length extrapolation, and determine model performance on the validation dataset. Additionally, we analyze their results on a set of downstream tasks, gauge their language modeling effectiveness across 571 domains from the PALOMA benchmark, and examine their properties with respect to scaling laws.

\paragraph{Sequence Length Extrapolation.}
\label{sec:spaj300B_ctx_extrapolation}
We begin by evaluating the ability of large-scale, 1.3B-parameter xLSTM, RWKV-4, Llama, and Mamba models to extrapolate to longer sequence lengths. Each model is trained using a context length of 2048, but at test time, we extend the context lengths up to 16384. The corresponding results are shown in Figure~\ref{fig:spaj300B_seq_extrapolation}. Notably, xLSTM demonstrates a consistent ability to keep perplexity low as the context length increases, distinguishing it from the other approaches.

\paragraph{Validation Perplexity and Downstream Tasks.}
\label{sec:spaj_downstream_eval}
Additionally, across every model scale, we assess how xLSTM, RWKV-4, Llama, and Mamba models perform on next token prediction using the SlimPajama validation dataset, as well as on downstream evaluations that probe common sense reasoning abilities.

The third column in Table~\ref{tab:lmeval} presents the perplexity scores on the validation set for the various approaches. For every model size, xLSTM[1:0] and xLSTM[7:1] exhibit the lowest validation set perplexities, indicating their superiority among the tested models. The remaining columns in Table~\ref{tab:lmeval} summarize results for a range of downstream tasks. In most cases and for all evaluated model sizes, xLSTM achieves the highest performance, with Mamba surpassing it only on the ARC task in select instances. Further information can be found in Appendix~\ref{sec:appExpLanguage}.

\paragraph{Performance on PALOMA Language Tasks.} Thirdly, we evaluate models of various sizes—including xLSTM, RWKV-4, Llama, and Mamba—on their ability to predict the next token across PALOMA language tasks~\citep{Magnusson:23arxivshort}. The metric for assessment is perplexity in next token prediction, spanning 571 different text domains that encompass sources such as nytimes.com and Reddit’s r/depression.
The results in Table~\ref{tab:ppleval} present perplexity scores, with language modeling domains covered in the first seven columns and domain-specific benchmarks summarized in the final five columns. On these language tasks, xLSTM[1:0] yields lower perplexity—indicating better performance—compared to xLSTM[7:1].

xLSTM[1:0] demonstrates reduced perplexity compared to Mamba in 568 of 571 text domains (99.5\%), surpasses Llama in 486 domains (85.1\%) with lower perplexity, and outperforms RWKV-4 in 570 of 571 domains (99.8\%) with respect to perplexity. Additional information is provided in Appendix~\ref{sec:appExpLanguage}.

\paragraph{Scaling Laws.} Fourth, we investigate the manifestation of power-law scaling, which provides a framework for projecting model performance at increased parameter counts~\citep{Kaplan:20,Brown:20short}. Figure~\ref{fig:temp_sclaw300B} illustrates the scaling trajectories of the models under study. While the scaling trends are qualitatively alike across all models, their baseline performances differ. Among the evaluated architectures, RWKV-4 demonstrates the lowest performance, with Llama and Mamba yielding intermediate results. Notably, xLSTM outperforms Mamba, maintaining a performance differential comparable in size to the gap between Mamba and Llama. These scaling patterns suggest that, as model size increases, xLSTM is likely to maintain its advantage over both Transformer-based and State-Space-based counterparts.

\textbf{Generation Times and Maximal Throughput.} We report measurements for text generation latency in Figure~\ref{fig:inference_time_generation} and for peak throughput in Figure~\ref{fig:inference_time_decoding_throughput} (left), focusing on our xLSTM models at the 1.3B parameter level. As baselines, we use equivalently sized variants of Mamba, Llama, and RWKV sourced from HuggingFace; for Llama, a static key-value cache is utilized. During experimentation, attempts to compile the complete Transformer cache and to employ \texttt{torch.compile} for Mamba were unsuccessful.

All models involved in the generation benchmarks were evaluated with a batch size of 1 and a pre-fill of 16 tokens, intentionally selected to provide the Transformer model with the most advantageous setup. In Figure~\ref{fig:inference_time_generation}, both xLSTM and the recurrent baselines Mamba and RWKV-4 exhibit linear complexity with respect to sequence length, contrasting with the quadratic cost seen for Llama. For the decoding throughput analyses, we experimented with varying batch sizes and pre-fills for Llama.

From Figure~\ref{fig:inference_time_decoding_throughput} (right), it is evident that xLSTM supports substantially higher batch sizes than Llama, attributable to its memory usage remaining constant. As a result, xLSTM attains the leading throughput figures among the models tested.

\section{Limitations}

(i) Unlike mLSTM, the design of sLSTM enforces sequential memory mixing, restricting opportunities for parallel computation and making efficient parallel execution unattainable. Nonetheless, we implemented an optimized CUDA kernel for sLSTM that achieves runtime performance within a factor of two relative to our parallelized mLSTM kernel. (ii) The mLSTM CUDA kernels currently lack advanced optimizations and, as such, exhibit execution speeds approximately four times slower than FlashAttention or the scan approach utilized in Mamba. Leveraging techniques similar to those in FlashAttention should enable the development of more efficient CUDA kernels. (iii) The matrix-based memory in mLSTM imposes considerable computational demand because it processes matrices of size $d \times d$. However, since memory update and access operations are parameter-free and can be executed in parallel via standard matrix routines, the added computation time from handling complex memory structures remains relatively modest. (iv) It is essential to pay close attention to how forget gates are initialized. (v) Since matrix memory size does not scale with sequence length, increasing the context window could potentially saturate the memory in the case of very long sequences. However, up to sequence lengths of 16k, this issue has not manifested as problematic (see Section~\ref{sec:spaj300B_ctx_extrapolation}). (vi) Because of the substantial computational resources required for large-scale language modeling, we have not conducted comprehensive architecture or hyperparameter optimization, particularly for large xLSTM models. Thorough optimization remains necessary for xLSTM architectures to fully realize their capabilities.

\section{Conclusion}
We have addressed, in part, our initial inquiry: To what extent can language modeling be advanced by scaling LSTM architectures to the order of billions of parameters? Our current findings suggest that such scaling allows performance to reach levels comparable with modern approaches, including Transformers and State Space Models. By introducing exponential gating with memory mixing and revising the memory framework, we have transformed the conventional LSTM into xLSTM. Evaluations demonstrate that xLSTM achieves competitive results in language modeling tasks when benchmarked against leading models such as Transformers and State Space Models. Analysis of scaling laws further suggests that as xLSTM models increase in size, they are poised to emerge as strong alternatives to the prevailing Large Language Models that utilize Transformer-based architectures. Additionally, xLSTM holds significant promise for broader applications across disciplines such as Reinforcement Learning, Time Series Prediction, and physical system modeling.

\end{document}